\documentclass[11pt,final,hidelinks]{article}
\makeatletter\def\input@path{{sty/}}\makeatother
\usepackage[nohints]{minitoc}
\usepackage{lmodern}
\usepackage[margin=1in]{geometry}
\usepackage[english]{babel}
\usepackage[utf8]{inputenc}
\usepackage[T1]{fontenc}
\usepackage[activate={true,nocompatibility},final,tracking=true,kerning=true,spacing=true,factor=1100,stretch=10,shrink=10]{microtype}
\microtypecontext{spacing=nonfrench}
\usepackage{graphicx}
\graphicspath{{img/}} 
\usepackage{algorithm2e}
\usepackage{caption}
\usepackage{bbm}
\usepackage{booktabs}
\usepackage{mathtools}
\usepackage{subfiles}
\usepackage{commath}
\usepackage{amsmath}
\usepackage{amsthm}
\usepackage{enumerate}
\usepackage{amssymb}
\usepackage{pgfplots}
\usepackage{tabu}
\usepackage{numprint}
\usepackage{hyperref}
\usepackage{tikz}
\usepackage{tikzscale}
\usepackage[square,numbers]{natbib}
\usepackage{cleveref}
\usepackage{siunitx}
\usepackage[section]{placeins}
\usepackage{listings}
\usepackage{minted}
\usepackage{backnaur}
\usepackage[symbols,toc,acronym]{glossaries}

\newcommand{\COMMENT}[1]{}
\makeglossaries
\loadglsentries{gls-custom}
\loadglsentries{gls-approxset}
\loadglsentries{gls-set}
\loadglsentries{gls-hashing}

%% Paper-specific definitions for the Bernoulli Hash Function paper.
% General notation is provided by alex.sty.

%% Probability mass function (not in alex.sty)
\NewDocumentCommand{\PDF}{mO{}}{%
    \ifx\\#2\\%
        \operatorname{f}\!\left(#1\right)%
    \else%
        \operatorname{f}_{#2}\!\left(#1\right)%
    \fi%
}

%% Concatenation operator (alias for \cat from alex.sty, used in text)
\newcommand{\catop}{\cat}

%% Bit string notation
\newcommand{\cisb}{\BitSet^{*}}
\newcommand{\cisbn}[1]{\BitSet^{#1}}

%% Indicator function
\newcommand{\indicator}[1]{\mathbbm{1}_{#1}}

%% Ordinal helpers
\newcommand{\ith}{$i^{\text{th}}$}
\newcommand{\jth}{$j^{\text{th}}$}
\newcommand{\kth}{$k^{\text{th}}$}
\newcommand{\nth}{$n^{\text{th}}$}

%% Algorithm2e function keywords
\SetKwFunction{MakeBHF}{MakeBHF}
\SetKwFunction{Contains}{Contains}
\SetKwFunction{HasKey}{HasKey}
\SetKwFunction{Encode}{encode}
\SetKwFunction{Decode}{decode}
\SetKwFunction{ro}{$h^*$}
\SetKwFunction{sampler}{BitLengthSampler}
\SetKwFunction{Cardinality}{cardinality}
\SetKwFunction{Count}{count}
\SetKwFunction{FalsePositiveRate}{false\_positive\_rate}

%% Algorithm2e data keywords
\SetKwData{collisions}{collisions}

%% Algorithm styling
\definecolor{ltgray}{rgb}{0.6,0.6,0.6}
\newcommand\mycommfont[1]{\footnotesize\ttfamily\textcolor{ltgray}{#1}}
\SetCommentSty{mycommfont}

%% TikZ libraries
\usetikzlibrary{arrows,shapes,positioning}
\usetikzlibrary{calc,decorations.markings}

%% Graphics path
\graphicspath{{img/}}

%% Null value for maps
\newcommand{\nullvalue}{\bot}

%% Adaptive threshold macros
\newcommand{\betadist}{\operatorname{Beta}}
\newcommand{\OS}[2]{R_{(#1:#2)}}   % p-th order statistic of m values

%% Algorithm2e keyword for adaptive construction
\SetKwFunction{MakeBHFAdaptive}{MakeBHF-Adaptive}


\usepackage{functionnotation}
\usepackage{probabilitynotation}
\usepackage{algorithmnotation}
\usepackage{matrixnotation}
\usepackage[fancy]{setnotation}
\usepackage[fancy]{relationnotation}
\usepackage{approxsetnotation}
\usepackage[section]{envnotation}

% Paper-specific macros (PHF domain)
\newcommand{\PHF}{\ensuremath{\operatorname{PHF}}}
\newcommand{\PH}{\ensuremath{\operatorname{PH}}}
\renewcommand{\ph}{\ensuremath{\operatorname{ph}}}
\newcommand{\phf}{\ensuremath{\operatorname{phf}}}
\newcommand{\ro}{\ensuremath{\operatorname{h^*}}}
\newcommand{\trunc}{\ensuremath{\operatorname{trunc}}}
\renewcommand{\BL}{\ensuremath{\operatorname{BL}}}
\renewcommand{\RE}{\ensuremath{\operatorname{e}}}
\renewcommand{\AE}{\ensuremath{\operatorname{e^*}}}
\newcommand{\BitMatrix}{\ensuremath{\operatorname{BitMatrix}}}
\newcommand{\pmapf}{\ensuremath{\operatorname{pmf}}}

\PassOptionsToPackage{unicode}{hyperref}
\PassOptionsToPackage{naturalnames}{hyperref}

\definecolor{ltgray}{rgb}{0.7,0.7,0.7}

\SetKwFunction{Encode}{encode}
\SetKwFunction{Rows}{nrows}
\SetKwFunction{Columns}{ncols}
\SetKwFunction{Contains}{contains}
\SetKwFunction{MakeApproxSet}{make\_approximate\_set}
\SetKwFunction{MakeApproxMap}{make\_approximate\_map}
\SetKwFunction{MakeMinimalPerfectHash}{make\_perfect\_hash}
\SetKwFunction{MakeKPerfectHashSet}{make\_k\_perfect\_hash\_filter}
\SetKwFunction{MakeTypeApproxSet}{make\_typed\_approximate\_set}

% Map and application macros
\newcommand{\Mt}{\mathcal{M}}
\newcommand{\notfn}{\operatorname{not}}
\newcommand{\andfn}{\operatorname{and}}
\newcommand{\orfn}{\operatorname{or}}
\newcommand{\diffn}{\operatorname{diff}}
\newcommand{\Has}{\operatorname{has}}
\newcommand{\truefn}{\operatorname{true}}
\newcommand{\very}{\operatorname{very}}


\lstset
{
	basicstyle=\footnotesize,               % print whole listing small
	keywordstyle=\bfseries,
	identifierstyle=,               % nothing happens
	commentstyle=\itshape\color{ltgray},    % gray comments
	%stringstyle=\ttfamily,          % typewriter type for strings
	showstringspaces=false,         % no special string spaces
	numbers=left,
	numberstyle=\tiny,
	stepnumber=1,
	numbersep=5pt
}

\hypersetup{
	pdftitle={The \emph{Perfect Hash Filter}},  % title
	pdfauthor={Alexander Towell},               % author
	pdfsubject={computer science},              % subject of the %document
	pdfkeywords={
		probabilistic data structure,
		abstract data type,
		Bernoulli set,
		Bernoulli map,
		bloom filter,
		perfect hash function,
		perfect hash filter,
		perfect map filter},                  % keywords
	colorlinks=true,                           % false: boxed links;
	linkcolor=blue,
	citecolor=green,                            % color of links to 
	filecolor=magenta,                          % color of file links
	urlcolor=green                              % color of external
}


\title
{
    The \emph{\gls*{gls-phf}}\\
    \large An efficient implementation of the positive Bernoulli set and Bernoulli map abstract data types
}
\author
{
    Alexander Towell\\
    \texttt{atowell@siue.edu}
}
\date{}

\begin{document}
\maketitle
\begin{abstract}
The \gls*{gls-phf} places each member of a static set at a unique slot via a minimal perfect hash and stores a short fingerprint there, the well-known static approximate-membership-query construction shared by Bloomier, xor, and cuckoo filters; the specific pairing of a minimal perfect hash with $\lceil\log_2(1/\varepsilon)\rceil$-bit fingerprints, and its $\log_2 e$ bits-per-element overhead, is already described by Dietzfelbinger and Pagh~\cite{dietzfelbinger2008}. The space-optimal construction of the Bernoulli set and map is the Bernoulli hash function~\cite{shs}, which reaches the information-theoretic floor $-\log_2 \varepsilon + \mu$ up to lower-order terms by a single-seed search whose expected construction work is exponential in the number of keys~\cite{lehmann2025seeds}; the \gls*{gls-phf} is instead the \emph{practical} construction, paying a constant $\log_2 e \approx 1.44$ bits per element above that floor in exchange for fast, deployable construction via off-the-shelf perfect hashing and a cheap tunable false negative rate. No novelty is claimed for the base structure or its analysis; the paper restates it within Bernoulli type theory, in four parts. First, we realize the \gls*{gls-phf} as a faithful implementation of the static \emph{Bernoulli map}~\cite{bmaps} (the \gls*{gls-pmf}), with the static \emph{positive Bernoulli set}~\cite{bsets} (false positive rate $\varepsilon = 2^{-M}$, zero false negatives) as the special case where the stored value is the membership indicator (zero-length value, $\mu = 0$), the same set-as-map unification the Bernoulli hash function uses; under the Bernoulli model it satisfies element-wise independence and conditional independence of rates and inherits the framework's error model and composition theory. Second, replacing the perfect hash with a \emph{$k$-perfect} hash~\cite{chd} admits a nonzero false negative rate that is cheap to construct, giving a two-rate structure whose false negatives are dependent within shared slots; it is not space-efficient, since dropping members outright, as in lossy dictionaries~\cite{pagh2001lossy}, uses fewer bits per retained member. Third, we develop the \gls*{gls-pmf} for general values (the map form is the Bloomier filter and the retrieval data structure~\cite{bloomier,dietzfelbinger2008}) and apply it to fuzzy set queries for rank-ordered retrieval. Fourth, we tie the construction's space to the information-theoretic lower bound of $-\log_2 \varepsilon$~bits per element due to Carter et al.~\cite{carter1978}: the \gls*{gls-phf}'s $\log_2 e$ overhead is a vanishing \emph{fraction} of the optimum as $\varepsilon \to 0$, so it is succinct in the relative sense, against the Bloom filter's $1.44$ multiplicative overhead, while we concede the absolute optimum to the Bernoulli hash function.
\end{abstract}
% Status and errata box (October 2026 freeze). \input right after \end{abstract}.
\begingroup
\setlength{\fboxsep}{7pt}%
\begin{center}
\fbox{\begin{minipage}{0.93\linewidth}\small
\setlength{\parskip}{4pt}%
\setlength{\parindent}{0pt}%
\textbf{Status and errata (October 2026).}
This is the final version of this preprint; no further revisions are planned.
It is superseded by Chapter~15 (\emph{Perfect Hash Filters}) of the monograph
\emph{Bernoulli Types: Two-Sided Error in Approximate Membership, Retrieval,
and Computation} (A.~Towell, 2026),
\url{https://metafunctor.com/latex/bernoulli-monograph/book.html}
(project hub: \url{https://metafunctor.com/research/bernoulli/}).
Please cite the monograph in preference to this preprint.

\textbf{Corrections in this version.}
\begin{itemize}\setlength{\itemsep}{1pt}\setlength{\parskip}{0pt}
\item Abstract, introduction, related work: the construction is credited to Dietzfelbinger and Pagh (2008, Sec.~1); novelty claims are withdrawn.
\item Introduction, related work: removed the claim that salt search needs combinatorial subset enumeration to admit false negatives; salt search is single-seed search with exponential expected work.
\item Sec.~5: the false negative formula ignored the rejection of overfull seeds; the exact conditioned law is added ($m=N=2$, $k=1$: true rate $0$, formula $1/4$).
\item Sec.~5: a displaced member whose fingerprint matches the retained one is still accepted; the false negative rate is $(1-2^{-M})\,\omega$ ($1/8$, not $1/4$, for $m=N=2$, $k=2$, $M=1$).
\item Sec.~5: the false negative count is not binomial (losses in a shared slot are dependent), so the $k$-perfect variant is not an element-wise independent Bernoulli set.
\item Sec.~3 footnote and Sec.~5: reversed the claim that $k$-perfect hashing is the more space-efficient way to admit false negatives; a lossy subset uses fewer bits per retained member.
\item Sec.~3 and Sec.~6: the perfect-hash seed is a geometric index, not a string of fair coins; the maximum-entropy statements now hold to leading order only.
\item Sec.~3: the space bound with false negatives is credited to Pagh and Rodler (2001, Prop.~1).
\end{itemize}

\textbf{Predecessors credited.} Minimal perfect hash plus fingerprints:
Dietzfelbinger and Pagh~\cite{dietzfelbinger2008}; static maps: Bloomier
filters~\cite{bloomier}; $k$-perfect hashing: Belazzougui, Botelho, and
Dietzfelbinger~\cite{chd}; false negatives traded for space: Pagh and
Rodler~\cite{pagh2001lossy}; the space bound: Carter et al.~\cite{carter1978};
seed search: Lehmann et al.~\cite{lehmann2025seeds}. The analyses here are
standard and no novelty is claimed for them.
A dated list of changes is in \texttt{ERRATA.md} in the source repository.
\end{minipage}}
\end{center}
\endgroup


\tableofcontents
%\listoffigures
%\listofalgorithms

\subfile{sections/intro}
\subfile{sections/related_work}
\subfile{sections/rd_ph_filter}
\subfile{sections/typed_phf}
\subfile{sections/k_perfect_hashing}
\subfile{sections/pmf}
\subfile{sections/apps_postings_list}
\subfile{sections/implementation}
\subfile{sections/adt}
\printglossary
\bibliographystyle{plainnat}
\bibliography{references}
\end{document}
